
\SGARefTarget{sga3:xi:chapter:xi}\chapter*{Exposé XI\\
Representability Criteria\\
Applications to Multiplicative-Type Subgroups\\
of Affine Group Schemes}
\addcontentsline{toc}{chapter}{Exposé XI. Representability Criteria and
Multiplicative-Type Subgroups}
\markboth{EXPOSÉ XI. REPRESENTABILITY CRITERIA}
  {EXPOSÉ XI. REPRESENTABILITY CRITERIA}

\begin{center}
  \textit{by A. Grothendieck}
\end{center}

\SGARefTarget{sga3:xi:section:0}\section*{0. Introduction}
\addcontentsline{toc}{section}{0. Introduction}

As we have already seen in Exposé X, nos.~4 and 5, the representability
of certain functors---for example, functors of the form
\(\operatorname{Hom}_S(X,Y)\) and various related constructions---plays
an important role in many questions concerning group preschemes.

Among the results that are particularly useful in this theory, besides
the representability of quotients studied in Exposés V and
\(\mathrm{VI}_{\mathrm B}\) and in Exposé VIII, no.~5, let us mention the
representability of functors of the form
\[
  \prod_{X/S} Y/X ,
\]
where \(Y\) is a subobject of \(X\).  This was studied in a very
elementary case in Exposé VIII, no.~6; variants will be given in no.~\SGARefLink{sga3:xi:section:6}{6}
of the present exposé.  These results provide the representability of
various centralizers, normalizers, and transporters.

Less elementary representability criteria, using results that will
appear in EGA VI, are indicated in \SGARefLink{sga3:xi:remark:6:12}{6.12} and in Exposés XV and XVI.
There we shall give a criterion for the representability of quotients
\(G/H\) in cases not covered by the preceding exposés; this criterion
was not developed in the oral lectures.

Our principal purpose in the present exposé is to prove Theorems \SGARefLink{sga3:xi:theorem:4:1}{4.1}
and \SGARefLink{sga3:xi:corollary:4:2}{4.2}, which provide a typical example of a nonprojective
construction technique, close to the technique that will be developed
in EGA VI.  Since the oral lecture and the drafting of the present
text, it has moreover become apparent that the affine hypotheses in
\SGARefLink{sga3:xi:theorem:4:1}{4.1} and \SGARefLink{sga3:xi:corollary:4:2}{4.2} can be removed to a large extent (cf.\ Exposé XV), and
that the main part of the theory developed in the following exposé can
dispense with \SGARefLink{sga3:xi:theorem:4:1}{4.1} and \SGARefLink{sga3:xi:corollary:4:2}{4.2}.  Finally, it might be interesting to prove
analogues of these results for a general reductive group prescheme
(for example, a semisimple one) in place of a group of multiplicative
type; in that case \SGARefLink{sga3:xi:theorem:4:1}{4.1} and \SGARefLink{sga3:xi:corollary:4:2}{4.2} would doubtless be the key result in the
proof.

\SGARefTarget{sga3:xi:section:1}\section*{1. Review of Smooth, Étale,
and Unramified Morphisms}
\addcontentsline{toc}{section}{1. Review of Smooth, Étale, and
Unramified Morphisms}

The reader is referred to EGA IV, §§17 and 18, and, pending its
publication, to SGA 1, Exposés I, II, and III.  Certain noetherian
hypotheses used there, which are inconvenient in applications, should
however be replaced by finite-presentation hypotheses.

\medskip
\noindent\SGARefTarget{sga3:xi:definition:1:1}\textbf{Definition 1.1.} ---
\textit{Let \(S\) be a prescheme and let
\[
  F:(\mathbf{Sch})^\circ/S\longrightarrow(\mathbf{Ens})
\]
be a functor.  The functor \(F\) is said to be \emph{formally smooth}
(respectively, \emph{formally unramified},
respectively, \emph{formally étale}) if, for every \(S\)-prescheme
\(S'\) that is affine in the absolute sense and every subscheme
\(S'_0\subset S'\) defined by a nilpotent ideal \(\mathscr J\), the map
\[
  F(S')\longrightarrow F(S'_0)
\]
is surjective (respectively, injective, respectively, bijective).}

\textit{An \(S\)-prescheme \(X\) is said to be formally smooth over
\(S\) (respectively, formally unramified over \(S\), respectively,
formally étale over \(S\)) if the corresponding functor has that
property.  It is said to be smooth over \(S\) (respectively,
unramified over \(S\), respectively, étale over \(S\)) if it satisfies
the preceding condition and is, in addition, locally of finite
presentation over \(S\).}

The usefulness of these definitions for \(X\) lies, on the one hand,
in their remarkably simple expression in terms of the functor
represented by \(X\) (and in practice \(X\) is often given as the
object over \(S\) representing an explicitly described functor), and,
on the other hand, in the remarkable local structural properties that
they express.  We recall those properties in the following statements;
for their proofs, see the references just cited.

\medskip
\noindent\SGARefTarget{sga3:xi:proposition:1:2}\textbf{Proposition 1.2.} ---
\textit{Let \(X\) be a prescheme locally of finite presentation over
\(S\).  Then:}

\begin{enumerate}
\item[(i)] \SGARefTarget{sga3:xi:proposition:1:2:item:i} \textit{\(X\) is smooth over \(S\) if and only if it is flat
over \(S\) and its geometric fibers
\[
  X\otimes_S\operatorname{Spec}\kappa(s)
\]
are regular schemes.  More generally, \(X\) is smooth over \(S\) in a
neighborhood of \(x\in X\) (one then says that \(X\) is smooth over
\(S\) at \(x\)) if and only if \(X\) is flat over \(S\) at \(x\) and
\(X_s\) is smooth over \(\kappa(s)\) at \(x\); that is,
\[
  X_s\otimes_{\kappa(s)}\overline{\kappa(s)}
\]
is regular at the points (or merely at one point) lying over \(x\).}

\item[(ii)] \SGARefTarget{sga3:xi:proposition:1:2:item:ii} \textit{Suppose that \(S\), and hence \(X\), is locally
noetherian.  Let \(x\in X\), and let \(s\in S\) be its image.  The
smoothness of \(X\) over \(S\) at \(x\) can then be detected from the
local homomorphism
\[
  A=\mathscr O_{S,s}\longrightarrow B=\mathscr O_{X,x}
\]
of noetherian local rings (or even from the local homomorphism
\(\widehat A\to\widehat B\) of their completions), by the following
characteristic property: \(B\) is flat over \(A\), and
\(B\otimes_A k\) is geometrically regular over \(k\), the residue
field of \(A\).  Equivalently, for every finite extension \(k'\) of
\(k\),
\[
  (B\otimes_A k)\otimes_k k'=B\otimes_A k'
\]
is a regular semilocal ring.  When the residue-field extension
\(k(B)/k(A)\) is trivial, these conditions are also equivalent to the
following one: \(\widehat B\) is isomorphic, as an
\(\widehat A\)-algebra, to an algebra of formal power series
\[
  \widehat A[[t_1,\ldots,t_n]].
\]}
\end{enumerate}

Thus, from the ``formal'' point of view, the structure of \(X\) over
\(S\) is that of the standard affine space \(S[t_1,\ldots,t_n]\) over
\(S\).

\medskip
\noindent\SGARefTarget{sga3:xi:proposition:1:3}\textbf{Proposition 1.3.} ---
\textit{Let \(X\) be a prescheme locally of finite presentation over
\(S\).  Then:}

\begin{enumerate}
\item[(i)] \SGARefTarget{sga3:xi:proposition:1:3:item:i} \textit{The following conditions are equivalent:}
  \begin{enumerate}
  \item[(a)] \SGARefTarget{sga3:xi:proposition:1:3:item:i:item:a} \textit{\(X\) is unramified over \(S\).}
  \item[(b)] \SGARefTarget{sga3:xi:proposition:1:3:item:i:item:b} \textit{The diagonal morphism
  \(X\to X\times_S X\) is an open immersion.}
  \item[(c)] \SGARefTarget{sga3:xi:proposition:1:3:item:i:item:c} \textit{\(\Omega^1_{X/S}=0\).}
  \item[(d)] \SGARefTarget{sga3:xi:proposition:1:3:item:i:item:d} \textit{The geometric fibers of \(X/S\) are discrete and
  reduced; that is, they are isomorphic to disjoint sums of copies of
  the base field.}
  \item[(e)] \SGARefTarget{sga3:xi:proposition:1:3:item:i:item:e} \textit{For every \(s\in S\), the fiber
  \[
    X\otimes_S\operatorname{Spec}\kappa(s)=X_s
  \]
  is unramified over \(\kappa(s)\), or equivalently is isomorphic to a
  disjoint sum of spectra of finite separable extensions of
  \(\kappa(s)\).}
  \end{enumerate}

\item[(ii)] \SGARefTarget{sga3:xi:proposition:1:3:item:ii} \textit{There are analogous pointwise conditions for
\(X\) to be unramified over \(S\) at a given point \(x\), that is, in
a neighborhood of that point.  For example, it is necessary and
sufficient that \(\mathscr O_{X_s,x}\) be a finite separable extension
of \(\kappa(s)\).  In terms of the local homomorphism
\[
  A=\mathscr O_{S,s}\longrightarrow B=\mathscr O_{X,x},
\]
this says that \(B\otimes_A k\) is a finite separable extension of
the residue field \(k\) of \(A\); equivalently,
\[
  \mathfrak r(A)B=\mathfrak r(B)
  \quad\text{and}\quad
  k(B)/k(A)\ \text{is finite separable}.
\]
When \(A\), and hence \(B\), is noetherian and
\(k(A)\xrightarrow{\sim}k(B)\), this also means that
\(\widehat A\to\widehat B\) is surjective.}
\end{enumerate}

Thus, from the ``formal'' point of view, saying that \(X\) is
unramified over \(S\) means that \(X\) is essentially a subprescheme
of \(S\).  Observe also that conditions \SGARefLink{sga3:xi:proposition:1:3:item:i:item:d}{(d)} and \SGARefLink{sga3:xi:proposition:1:3:item:i:item:e}{(e)} are expressed
solely in terms of the fibers of \(X/S\).

\medskip
\noindent\SGARefTarget{sga3:xi:proposition:1:4}\textbf{Proposition 1.4.} ---
\textit{Let \(X\) be a prescheme locally of finite presentation over
\(S\).  Then \(X\) is étale over \(S\) if and only if it is smooth and
unramified over \(S\) (which is immediate from the definition), so the
criteria in \SGARefLink{sga3:xi:proposition:1:2}{1.2} and \SGARefLink{sga3:xi:proposition:1:3}{1.3} apply.  In particular:}

\begin{enumerate}
\item[(i)] \SGARefTarget{sga3:xi:proposition:1:4:item:i} \textit{\(X\) is étale over \(S\) if and only if it is flat
and unramified over \(S\).  There is an analogous local criterion for
\(X\) to be étale over \(S\) at a given point \(x\).}

\item[(ii)] \SGARefTarget{sga3:xi:proposition:1:4:item:ii} \textit{Suppose in addition that \(S\), and hence \(X\),
is locally noetherian.  Whether \(X\) is étale over \(S\) at \(x\) can
then be detected from the local homomorphism
\[
  A=\mathscr O_{S,s}\longrightarrow B=\mathscr O_{X,x}
\]
\footnote{The source prints \(\mathscr O_{X,s}\); the local ring at the
point of \(X\) is \(\mathscr O_{X,x}\).}
(and even from \(\widehat A\to\widehat B\)) by the following
characteristic property: \(B\) is flat over \(A\), and \(B\otimes_A k\)
is a finite separable extension of \(k\), the residue
field of \(A\).  When
\(k(A)\xrightarrow{\sim}k(B)\), this condition simply says that
\(\widehat A\to\widehat B\) is an isomorphism.}
\end{enumerate}

Thus, from the ``formal'' point of view, saying that \(X\) is étale
over \(S\) simply means that \(X\) is locally isomorphic to \(S\).

\medskip
\noindent\SGARefTarget{sga3:xi:remark:1:5}\textbf{Remark 1.5.} ---
When \(X\) is locally of finite presentation over \(S\) and a point
\(x\in X\) is given, the fact that \(X\) is smooth (respectively,
unramified, respectively, étale) over \(S\) at \(x\), that is, in a
neighborhood of \(x\), can also be detected from the functor
\[
  X:(\mathbf{Sch})^\circ/S\longrightarrow(\mathbf{Ens})
\]
by the following property.  For every \(S\)-scheme \(S'\) that is the
spectrum of a local ring, every subscheme \(S'_0\subset S'\) defined
by a nilpotent ideal, and every \(S\)-morphism
\(u_0:S'_0\to X\) sending the closed point \(s'\) of \(S'\) to \(x\),
there exists at least one (respectively, at most one, respectively,
exactly one) \(S\)-morphism \(u:S'\to X\) extending \(u_0\).

This statement shows in particular that, in \SGARefLink{sga3:xi:definition:1:1}{Definition 1.1}, one may
restrict to local schemes \(S'\), provided the functor under
consideration is representable by a prescheme locally of finite
presentation over \(S\).  Moreover, when \(S\) is locally noetherian,
one may even restrict the preceding pointwise criterion to local
artinian schemes \(S'\), and the same restriction may therefore be
made in \SGARefLink{sga3:xi:definition:1:1}{Definition 1.1}, provided that \(S\) is locally noetherian and
that the functor in question is represented by a prescheme locally of
finite type over \(S\).

